% =============================================================================
% Template thesis using the qthesis class
% =============================================================================
% This file is the main entry point. Build: make (uses latexmk; see Makefile).
% Default engine: lualatex. To use another: make ENGINE=pdf or make ENGINE=xe.
%
% SPDX-FileCopyrightText: 2026 Frank Langbein <frank@langbein.org>
% SPDX-License-Identifier: LPPL-1.3c
% SPDX-License-Identifier: CC-BY-SA-4.0
%
% This template and other .tex/.bib files in this directory are part of the
% same work and may be distributed and/or modified under the LaTeX Project
% Public License (LPPL), version 1.3c. See LICENSE file.
%
% -----------------------------------------------------------------------------
% Overview (qthesis class)
% -----------------------------------------------------------------------------
% The class is a wrapper around the standard book class. Layout: A4, 12pt,
% 3cm margins; 1.5 line spacing; ragged-right; main text 12pt, footnotes and
% captions at least 11pt; preliminary pages in roman numerals, then Arabic
% for main matter. Chapter style: top bar in header, ``Chapter N'' or
% ``Appendix A'' in header; number/letter and title on first line of chapter.
%
% -----------------------------------------------------------------------------
% Thesis structure
% -----------------------------------------------------------------------------
% Main body: Introduction, Literature Review, Methods, Results, Discussion,
% optional Further Analysis, Conclusions and Future Work. Adapt as needed
% (e.g.\@ multiple results or theory chapters, merged Discussion/Results).
% See each chapter file for suggested sections and alternatives.
%
% Suggested sections at a glance (omit/rename as needed):
%   1 Introduction:    Motivation and context; Research questions and objectives;
%                      Contributions; Thesis structure (roadmap). Optional: Scope, Terminology.
%   2 Literature Review: Thematic sections (e.g.\@ by area or chronology); end with
%                      gap and positioning. Optional: Theory and concepts; subsections by theme.
%   3 Methods:         Research design; Design; Procedure; Implementation or experimental
%                      setup. Optional: Data/sources; Ethical considerations; Theory/formal framework.
%   4 Results:         Setup or experimental protocol; Main findings; optional Additional results.
%                      Optional subsections: by research question, by dataset or condition.
%   5 Discussion:      Interpretation; Limitations; Comparison with prior work.
%                      Optional: Implications; Threats to validity.
%   6 Further Analysis: Optional chapter: sensitivity analysis, extensions, case studies.
%                      Sections as needed (e.g.\@ Analysis, Implications).
%   7 Conclusions:     Conclusions (summary of aims and results; no new material); Future work.
%
% Structure variants (adapt as needed): thesis-by-publication (compilation of
% papers plus intro/conclusion chapters); multiple results or theory chapters
% (e.g.\@ Study 1, Study 2, or Theory, Experiment 1, Experiment 2); merged
% Results and Discussion. Ch6 (Further Analysis) can be omitted or folded into
% Results/Discussion. See comment in Ch2 for splitting Literature Review.
%
% -----------------------------------------------------------------------------
% Class options
% -----------------------------------------------------------------------------
% Font: sans (default), serif, or hacker;
%   serif = serif text and math.
% Layout: oneside (default) or twoside.
% Citation: numeric (default) = plainnat, [1], [1,2]; harvard = agsm, (Author, 2000).
% Alignment: raggedright (default) = ragged-right; justified = fully justified text.
% Line spacing: onehalfspacing (default); single; or linespread115/linespread12
%   for institutions allowing 1.0--1.3.
% Paragraph style: parskip (default) = space between paragraphs, no first-line indent;
%   parindent = traditional first-line indent, minimal parskip.
% Optional typography:
%   fullwidth  = enable \begin{fullwidth}...\end{fullwidth} and fullwidthfigure/
%                fullwidthtable environments that extend into the margin (requires
%                changepage or chngpage).
%   italic = chapter, section, subsection and captions in italic (default: bold/roman).
% Engine: pdfTeX vs XeLaTeX/LuaLaTeX is detected for font setup.
%
% Defaults (no option): link colour DarkBlue, cite/URL colour DarkGreen; tight list
% margins; quote, quotation, and verse in smaller type with indents;
% \epigraph{quote}{source} for chapter/section-opening quotes.
% See Section~\ref{sec:ch1-reference} in Chapter~1 for examples of \newthought,
% fullwidth, fullwidthfigure, quote, and epigraph.
%
% -----------------------------------------------------------------------------
% Metadata and title page
% -----------------------------------------------------------------------------
% \title, \author, \degree, \university, \date drive the generated title page
% (bold title, name, degree, university, month/year).
% Optional: \keywords{...} is typeset after the abstract (e.g.\@
% \keywords{keyword one; keyword two; keyword three}). Many institutions require keywords.
%
% -----------------------------------------------------------------------------
% License and copyright (SPDX)
% -----------------------------------------------------------------------------
% \copyrightholder defaults to \author; \licenseyear defaults to current year.
% When \license is set: (1) a copyright page is inserted after the abstract;
% (2) a License chapter is inserted just before the bibliography.
% If qthesis-licenses.tex is present, the class loads it and can show full
% license text for supported SPDX IDs; otherwise name + URL only.
% Supported SPDX IDs include: MIT, Apache-2.0, BSD-2/3-Clause, GPL/AGPL/GFDL,
% LPPL-1.3c, CC-*, CC0-1.0. Any SPDX id accepted; unlisted -> spdx.org.
%
% -----------------------------------------------------------------------------
% Front matter
% -----------------------------------------------------------------------------
% \maketitle{...} produces: title page; then (if \license set) copyright page;
% then abstract; then frontlists: TOC, LOF, LOT, LOA, List of Acronyms.
% Optional unnumbered chapters: \dedication{...}; Acknowledgements; Declaration;
% Publications, data and code. Use \chapter*{Title} and \addcontentsline{toc}{chapter}{Title}
% for those that are not provided by the class (e.g.\@ \dedication is a class command).
%
% -----------------------------------------------------------------------------
% Chapters and back matter
% -----------------------------------------------------------------------------
% \mainmatter starts Arabic page numbering. Put \backmatter before \bib.
% \bib{basename} runs bibliography (and inserts License chapter if \license set).
% Appendices: \appendix then \chapter for each; lettered A, B, ... in headers/TOC.
% Natbib: numeric -> plainnat; harvard -> agsm.
%
% -----------------------------------------------------------------------------
% Headings: use title case for chapter titles (capitalise major words; keep
% articles, conjunctions, and short prepositions lowercase unless first/last
% word). Use sentence case for sections, subsections, and below (only the
% first letter and proper nouns capitalised). Or adjust as preferred, but keep
% it consistent.
%
% -----------------------------------------------------------------------------
% LaTeX formatting (citations, spaces, dashes, quotes)
% -----------------------------------------------------------------------------
% Citations (natbib): prefer \citep and \citet over bare \cite (class may warn).
%   \citep{key}  = parenthetical: ``(Author, 2020)'' or ``[1]''. Use when the citation
%     is a clause or in parentheses: ``as shown~\citep{key}'' or ``(e.g.\@ \citep{key})''.
%   \citet{key}  = narrative: ``Author (2020)'' or ``Author [1]''. Use when the author
%     is part of the sentence. If the sentence *starts* with the citation, write
%     ``\citet{key} showed that...'' (no ~ before \citet at sentence start). If the
%     citation follows a word, use a non-breaking space: ``Smith~\citet{key} argued...''.
%   Non-breaking space (~): put ~ before \citep or \citet when they follow a word so
%     the citation does not line-break to the next line; omit ~ at sentence start.
% Spaces: \@ after abbreviation before period (e.g.\@ something); ~ non-breaking
%   (Fig.~1, Section~2.1); \, thin space (numbers, number+unit); "\ " normal
%   stretchable space; \nobreakspace or ~ to avoid line break.
% Dashes: -- en-dash (ranges, compound adjectives); --- em-dash (parenthetical).
% Ellipsis: \ldots. Quotes: use `` and '' (or csquotes) for typographic quotes.
% Reviewer comments: %ID (e.g.\@ %F) at line start, space after ID.
%
% -----------------------------------------------------------------------------
% Footnotes: prefer putting short asides in parentheses in the main text. Use
% footnotes only sparingly when they add value (e.g.\@ translations, optional
% clarifications that would interrupt the flow).
%
% -----------------------------------------------------------------------------
% Float placement (figures, tables, algorithms)
% -----------------------------------------------------------------------------
% Use the optional argument: \begin{figure}[placement] ... \end{figure}. Same for
% table, algorithm. Placement letters (LaTeX tries in order): h = here (approx),
% t = top of page, b = bottom of page, p = separate page of floats. Order matters:
%   [htbp]  try here, then top, then bottom, then float page (common default).
%   [t]     top only; good when you prefer figures/tables at page top.
%   [thbp]  prefer top, then here, bottom, float page.
% Prefix ! relaxes LaTeX's limits (e.g. allow more floats per page): use [!t],
% [!htbp] when you want stronger preference (e.g. \begin{figure}[!t] for ``try
% hard to put at top''). The float package adds [H] = ``here definitely'' (figure
% is not floated; use sparingly as it can cause uneven spacing).
% The class (qthesis.cls) relaxes float placement parameters so more floats fit
% on each page and full float pages are easier; floats tend to stay closer to
% their first reference in the text.
%
% -----------------------------------------------------------------------------
% Reviewer comments: convention ``%ID ''
% -----------------------------------------------------------------------------
% Format: %ID (e.g. %F for Frank) at start of line; space after ID required (grep-friendly).
% Purpose: comments from a person (ID): suggestions, questions on the *preceding* text.
% Placement: immediately after the sentence/element, on a new line (no break in paragraph).
% Find all: grep -r '^%[A-Za-z] ' *.tex */*.tex
% Lifecycle: when addressed, delete the comment; do not keep ``for checking''.
%
% -----------------------------------------------------------------------------
% Version control (git)
% -----------------------------------------------------------------------------
% Track sources (.tex, .bib, .cls, etc.) with git or another VCS.
% Ignore build products (see .gitignore in this directory for extensions).
% Use branches for major edits or chapter drafts; tag milestones (e.g. draft-v1, submission).
% Do not commit personal data or secrets; keep author, university, paths appropriate.
%
% -----------------------------------------------------------------------------
% Writing and project: Keep chapters focused (intro: motivation, questions, contributions,
% roadmap; literature: survey and gap; methods: reproducible; results/discussion/conclusions).
% Use \ref and \citep/\citet; label chapters/sections/figures; split into per-chapter .tex
% files; use version control and tag milestones; automate with make (or latexmk).
%
% -----------------------------------------------------------------------------
% Word count: run ``make wordcount'' (texcount -inc; per-file and total with caption
% words shown separately so you can apply institution rules if needed).
% =============================================================================
\documentclass[sans,oneside,numeric,raggedright,fullwidth]{qthesis}

% -----------------------------------------------------------------------------
% Optional: use a different acronyms file
% -----------------------------------------------------------------------------
%\renewcommand{\qthesisacronymfile}{my-acronyms.tex}

% -----------------------------------------------------------------------------
% Useful packages for thesis writing
% -----------------------------------------------------------------------------

% Tables
\usepackage{booktabs}        % \toprule, \midrule, \bottomrule (vertical rules generally not needed)
\usepackage{tabularray}      % Flexible column types, rules, longtblr for multi-page tables
\usepackage{longtable}       % Tables that span multiple pages (or use tabularray longtblr)
%\usepackage{multirow}        % Span multiple rows in tables
%\usepackage{multicol}        % Multiple columns in text/tables
%\usepackage{colortbl}        % Coloured rows/cells in tables

% Math, symbols, related
\usepackage{mathtools}       % Extensions to amsmath (matrix envs, \DeclarePairedDelimiter)
\usepackage{pifont}          % Dingbat symbols, e.g. \ding{51}, \ding{43}
\usepackage{siunitx}         % Units and numbers: \SI{3}{m/s}, \num{1000}

% Floats
\usepackage{float}           % [H] = here definitely; see ``Float placement'' comment block above
\usepackage{stfloats}        % Improved placement for bottom floats
\usepackage{subcaption}      % Subfigures with (a), (b) and per-subfigure captions
\usepackage{rotating}        % \begin{sidewaystable}, \begin{sidewaysfigure} for wide content

% Algorithms
\usepackage{algorithm}       % Numbered algorithm floats and List of Algorithms (default)
\usepackage[noend]{algpseudocode} % Algorithmicx pseudocode environment
% Alternative: \usepackage{algorithmic}  % Traditional algorithm layout, different command set
% Alternative: \usepackage[ruled,vlined]{algorithm2e}  % Different syntax and layout

% Other useful packages
\usepackage{lipsum}            % Placeholder text: \lipsum[n], \lipsum[1-5], etc.
\usepackage{enumitem}          % Customise list labels, spacing (e.g. [label=O\arabic*])
\usepackage{relsize}           % \textlarger, \textsmaller for relative font sizing
\usepackage[version=4]{mhchem} % Chemical formulae: \ce{H2O}, \ce{CO2}, \ce{^1H-NMR}
\usepackage{verbatim}          % \begin{verbatim} ... \end{verbatim} for code
\usepackage{afterpage}         % \afterpage{\clearpage} to clear after current page

% -----------------------------------------------------------------------------
% Thesis metadata
% -----------------------------------------------------------------------------
\title{Template Thesis Title: Each Word Capitalised}
\author{Your Full Name}
\degree{Doctor of Philosophy}
\university{Your University}
\date{February 2026}

% License and copyright following SPDX. (optional; comment out if not needed).
% Needs qthesis-licenses.tex for license text in the License chapter.
\license{CC-BY-SA-4.0}
\copyrightholder{Your Full Name}
\licenseyear{2026}

% Optional: keywords appear after the abstract.
\keywords{keyword one; keyword two; keyword three}

\begin{document}

\frontmatter

% \maketitle{...} produces: title page, (copyright page if \license set), abstract,
% then frontlists: TOC, List of Figures, List of Tables, List of Algorithms,
% List of Acronyms (if acronyms.tex defines \newacronym).
\maketitle{%
  This is the abstract: a short, self-contained summary of the thesis. It
  should state the problem, methods, and main contributions so that a reader
  can understand the scope without reading the full document. Optionally
  include one sentence on data/code availability (e.g.\@ ``Data and code are
  available as described in the chapter Publications, Data and Code''). The
  \texttt{qthesis} class typesets it as an unnumbered chapter and adds it to
  the table of contents. If you set \texttt{\string\keywords\{...\}}, they
  appear after the abstract.
}

% -----------------------------------------------------------------------------
% Optional front-matter chapters (unnumbered).
% -----------------------------------------------------------------------------

\dedication{%
  % Replace with your dedication (one or two lines); optional.
  To my family.
}

\chapter{Acknowledgements}
% Replace with your acknowledgements. Here you can also acknowledge any contributions from other
% parties to the work.
I thank my supervisors and colleagues for their support.

%\chapter{Declaration}
% Replace with your declaration (e.g. that the work is your own, and any
% required statements on collaboration, publications, or word count).
% Often this is submitted in separate forms and does not have to be here.
%This thesis is the result of my own work. Where reference is made to other work, this is indicated in the text.

\chapter{Publications, Data and Code}
% List publications and state where data and code produced for this thesis are available.
% Replace with your own entries; you can use \citep{key} for publications.

\section{Publications}
% Papers (and optionally other outputs) that form part of or arise from this thesis.
% These can also be discussed in the introduction (what has been published or is planned).
\begin{enumerate}
  \item Author. Title. \textit{Journal}, volume, pages, year.~\citep{example-paper}
\end{enumerate}

\section{Data}
% Datasets produced in this thesis and where they are available (repository, DOI, or ``on request'').
Data supporting this thesis are available from [repository/DOI or ``available on request'']. Replace with your statement.

\section{Code}
% Software or code produced in this thesis and where it is available (repository, DOI, etc.).
Code produced for this thesis is available from [repository/DOI or ``available on request'']. Replace with your statement.

\mainmatter

% -----------------------------------------------------------------------------
% Main chapters: one per folder C1--C7, each with chapterN.tex
% -----------------------------------------------------------------------------
\chapter{Introduction}\label{ch1}
% =============================================================================
% CHAPTER 1: Introduction (thesis template + feature demonstration)
% =============================================================================
% SPDX-FileCopyrightText: 2026 Frank Langbein <frank@langbein.org>
% SPDX-License-Identifier: LPPL-1.3c
%
% This chapter is structured as a typical thesis introduction: motivation,
% research questions, contributions, and thesis roadmap. Replace each section
% with your own content. It also demonstrates packages and class features
% (citations, figures, tables, etc.); see comments and Section~\ref{sec:ch1-reference}.
% Document-level conventions (citations, formatting, %ID comments) are in main.tex.
% =============================================================================

% -----------------------------------------------------------------------------
% Chapter intro (no section): one short paragraph before the first \section
% -----------------------------------------------------------------------------
This chapter introduces the thesis. It outlines the motivation and context for the work, the research questions and objectives, the main contributions, and the structure of the remaining chapters. The following sections expand on each of these.

% -----------------------------------------------------------------------------
% Motivation and context (typical first section)
% -----------------------------------------------------------------------------
% Use \section, \subsection, \label{key}. Refer with \ref{key} or \autoref{key}.

\section{Motivation and context}\label{sec:motivation}

This template illustrates the structure and style of a \gls{phd} thesis using the Qyber thesis class. As in a real introduction, you would state the problem domain, why it matters, and how your work fits. Acronyms use the long-short style: first use gives the full form (e.g.\@ \gls{uk}); later use \acrshort{phd} or \acrlong{phd} when you need only short or long.

As shown in previous work~\citep{example-paper}, we follow standard practice.
%F Clarify whether this applies to all cases.
\citet{example-book} provide a longer treatment; multiple citations: \citep{example-paper,example-conf}. You can link to other chapters (e.g.\@ Chapter~\ref{ch2}, Section~\ref{sec:objectives}) and to figures (Fig.~\ref{fig:demo}). The class provides \url{https://example.com} and \href{https://example.com}{link text} for URLs and hyperlinks.
%S Consider adding a reference here.

% -----------------------------------------------------------------------------
% Research questions and objectives (typical second section)
% -----------------------------------------------------------------------------
% Use \citep{} (parenthetical) or \citet{} (narrative), not \cite.

\section{Research questions and objectives}\label{sec:objectives}

State your main research question and break it down into objectives. The following is a placeholder list (enumitem package: custom labels).

\begin{enumerate}[label=\textbf{O\arabic*:}, leftmargin=*]
  \item First objective: replace with your own.
  \item Second objective.
  \item Third objective.
\end{enumerate}

When you need mathematics, use amsmath and mathtools. Example:
\begin{equation}\label{eq:example}
  E = mc^2, \qquad \int_0^1 f(x)\,\mathrm{d}x = F(1) - F(0).
\end{equation}
Refer as Equation~\eqref{eq:example}; inline math: $\alpha + \beta = \gamma$. Paired delimiters (mathtools): \DeclarePairedDelimiter{\abs}{\lvert}{\rvert} then $\abs{x}$, $\abs*{\frac{1}{2}}$. Chemistry (mhchem): \ce{H2O}, \ce{CO2}, \ce{^1H-NMR}. Units (siunitx): \SI{3.0}{m/s}, \num{1234} samples. Symbols: \textdegree, \textcopyright; dingbats (pifont): \ding{51} \ding{43}. Relative size (relsize): \textlarger{larger} \textsmaller{smaller}.

% -----------------------------------------------------------------------------
% Contributions (typical third section)
% -----------------------------------------------------------------------------

\section{Contributions}\label{sec:contributions}

Summarise the main contributions of the thesis. You can use a short list and refer to a figure or table.

\begin{itemize}
  \item First contribution; replace with your own.
  \item Second contribution.
\end{itemize}

Figure~\ref{fig:demo} shows how to include a figure (graphicx, caption; optional short caption for the List of Figures). Float placement: \texttt{[t]} = top of page (default here); \texttt{[htbp]} = here, top, bottom, float page; \texttt{[!t]} = relax limits and prefer top. See the ``Float placement'' comment block in \texttt{main.tex}. Use \texttt{[H]} (float package) only when you must force ``here'' (can cause uneven spacing). Subfigures (subcaption): see Figure~\ref{fig:sub}.

\begin{figure}[t]
  \centering
  \fbox{\rule{0pt}{4cm}\rule{4cm}{0pt}}
  \caption[Short for list of figures]{Full caption. Optional argument used in List of Figures.}\label{fig:demo}
\end{figure}

\begin{figure}[t]
  \centering
  \begin{subfigure}[t]{0.45\linewidth}
    \centering
    \fbox{\rule{0pt}{2cm}\rule{2cm}{0pt}}
    \subcaption{First subfigure.}
  \end{subfigure}
  \hfill
  \begin{subfigure}[t]{0.45\linewidth}
    \centering
    \fbox{\rule{0pt}{2cm}\rule{2cm}{0pt}}
    \subcaption{Second subfigure.}
  \end{subfigure}
  \caption{Example with subcaption (a) and (b).}
  \label{fig:sub}
\end{figure}

Tables: tabularray (default) with \texttt{tblr}, or longtable for multi-page tables (Table~\ref{tab:long}). Same placement options as figures (\texttt{[t]}, \texttt{[!t]}, \texttt{[htbp]}, etc.).

\begin{table}[t]
  \centering
  \caption{Example table (tabularray).}
  \label{tab:demo}
  \begin{tblr}{colspec={lcc}, hlines}
    Item   & Value A & Value B \\
    First  & 1.0     & 2.0    \\
    Second & 3.0     & 4.0    \\
  \end{tblr}
\end{table}

\begin{center}
  \begin{longtable}{lcc}
    \caption{Minimal longtable demo (not in a float environment).} \label{tab:long} \\
    \hline
    Item   & A & B \\ \hline
    \endfirsthead
    \hline
    Item   & A & B \\ \hline
    \endhead
    Row 1  & 1 & 2 \\
    Row 2  & 3 & 4 \\ \hline
  \end{longtable}
\end{center}

% Alternative (traditional): \usepackage{booktabs} then \toprule, \midrule, \bottomrule.

% -----------------------------------------------------------------------------
% Thesis structure / roadmap (typical final section of introduction)
% -----------------------------------------------------------------------------

\section{Thesis structure}\label{sec:roadmap}

Outline how the rest of the thesis is organised so the reader knows what to expect. Reference chapters by label.

\begin{itemize}
  \item \textbf{Chapter~\ref{ch2}} (Literature Review): surveys prior work and identifies the gap.
  \item \textbf{Chapter~\ref{ch3}} (Methods): describes design and procedure.
  \item \textbf{Chapter~\ref{ch4}} (Results): presents main findings.
  \item \textbf{Chapter~\ref{ch5}} (Discussion): interpretation and limitations.
  \item \textbf{Chapter~\ref{ch6}} (Further Analysis): extended or supplementary results.
  \item \textbf{Chapter~\ref{ch7}} (Conclusions): answers to research questions and future work.
  \item \textbf{Appendix~\ref{app:supp}}: supplementary material (tables, derivations, notation).
\end{itemize}

Algorithms (algorithm, algpseudocode) appear in the List of Algorithms (Algorithm~\ref{alg:demo}). Use \texttt{[t]} or \texttt{[!t]} for top placement like figures/tables. For code snippets use verbatim; for wide tables or figures use \texttt{\string
\begin{sidewaystable}} or \texttt{\string
  \begin{sidewaysfigure}} (rotating).\footnote{Footnote example. Captions and footnotes are set at least 11\,pt by the class.} Placeholder prose in other chapters uses \texttt{\string\lipsum[n]}. Example: \lipsum[1]
    Use \texttt{\string\afterpage\{\string\clearpage\}} (afterpage) when you need to clear after the current page.

    \begin{algorithm}[t]
      \caption{Example algorithm (algpseudocode). Alternative: algorithmic for traditional layout.}
      \label{alg:demo}
      \begin{algorithmic}[1]
        \State \textbf{Input:} data $x$
        \State \textbf{Output:} result $y$
        \State $y \gets 0$
        \For{$i = 1$ \textbf{to} $n$}
        \State $y \gets y + x_i$
        \EndFor
        \State \Return $y$
      \end{algorithmic}
    \end{algorithm}

    Short code (verbatim):
\begin{verbatim}
  x = 1
  y = x + 1
\end{verbatim}

    % -----------------------------------------------------------------------------
    % Template reference (for authors: where to find package demos and conventions)
    % -----------------------------------------------------------------------------

    \section{Template reference}\label{sec:ch1-reference}

    This introduction doubles as a template: the sections above (Motivation, Research questions, Contributions, Thesis structure) are the usual structure for a thesis Chapter~1. Replace their content with your own. The same file also demonstrates: \verb|\citep|/\verb|\citet| (natbib), glossaries (long-short acronyms), amsmath/mathtools, textcomp, pifont, siunitx, relsize, mhchem, enumitem, graphicx/caption, float/stfloats, subcaption, tabularray/longtable, algorithm/algpseudocode, verbatim, afterpage, lipsum, rotating, footnotes. Conventions for citations, spaces (\verb|~|, \verb|\@|), dashes, and reviewer comments (\verb|%ID |) are in the comment block at the top of \texttt{main.tex}.


\chapter{Literature Review}\label{ch2}
% =============================================================================
% CHAPTER 2: Literature Review
% =============================================================================
% SPDX-FileCopyrightText: 2026 Frank Langbein <frank@langbein.org>
% SPDX-License-Identifier: LPPL-1.3c
%
% Suggested sections: thematic sections (e.g.\@ by area or chronology); end with
% gap and positioning. Optional: Theory and concepts; subsections by theme.
% Alternatively split into multiple chapters (e.g.\@ Background, Related work) or
% by theme (Ch2a, Ch2b) if your committee prefers. Replace with your organised
% review and gap analysis.
% =============================================================================

This chapter surveys prior work relevant to the thesis and identifies the gap that this research addresses. Replace this paragraph and the lipsum blocks below with your organised review and gap analysis.

\lipsum[1-4]

\section{First theme}\label{sec:lit-first}
% By theme, area, or chronology; identify gap and position your work.

\lipsum[5-7]

\section{Second theme}\label{sec:lit-second}
% By theme, area, or chronology; end with gap and positioning.

\lipsum[8-10]


\chapter{Methods}\label{ch3}
% Placeholder chapter: replace with your methods.
% Lipsum: \lipsum (default 1--7), \lipsum[2], \lipsum[11-15], etc.
% Start with an intro paragraph (no \section) before the first section.
%
% SPDX-FileCopyrightText: 2026 Frank Langbein <frank@langbein.org>
% SPDX-License-Identifier: LPPL-1.3c

This chapter describes the methods used in the study. It covers design, procedure, and implementation. Replace with enough detail for reproduction.

\lipsum[11-14]

\section{Design}

\lipsum[15-17]

\section{Procedure}

\lipsum[18-20]

\section{Implementation}

\lipsum[66-67]


\chapter{Results}\label{ch4}
% =============================================================================
% CHAPTER 4: Results
% =============================================================================
% SPDX-FileCopyrightText: 2026 Frank Langbein <frank@langbein.org>
% SPDX-License-Identifier: LPPL-1.3c
%
% Suggested sections: Setup or experimental protocol; Main findings; optional
% Additional results. Optional subsections: by research question (e.g.\@ Results
% for RQ1, Results for RQ2), by dataset or condition. Structure by RQ and refer
% to \ref{rq1}, \ref{rq2} from Ch1. For statistical reporting, follow your
% discipline (e.g.\@ effect sizes, CIs, p-values; APA or field-specific standards).
% Replace with your figures, tables, and analysis.
% =============================================================================

This chapter presents the main results of the thesis. It is organised around the main findings and additional results. Replace with your figures, tables, and analysis.

\lipsum[23-24]

\section{Main findings}\label{sec:results-main}
% Optional subsections: by research question (e.g.\@ Results for RQ1), by dataset or condition.

\lipsum[25-27]

\paragraph{Key result.} A short paragraph here can highlight the main quantitative or qualitative finding before the reader continues to the full sections. Replace with your own key result or remove.

\section{Additional results}\label{sec:results-additional}

\lipsum[28-30]


\chapter{Discussion}\label{ch5}
% =============================================================================
% CHAPTER 5: Discussion
% =============================================================================
% SPDX-FileCopyrightText: 2026 Frank Langbein <frank@langbein.org>
% SPDX-License-Identifier: LPPL-1.3c
%
% Suggested sections: Interpretation; Limitations; Comparison with prior work.
% Optional: Answer to each research question (e.g.\@ interpretation and answer to RQ1);
% Theoretical and practical implications; Threats to validity.
% Replace with your discussion.
% =============================================================================

This chapter interprets the results, discusses limitations, and compares with prior work. Replace with your discussion.

\lipsum[31-34]

\section{Interpretation}\label{sec:disc-interp}

\lipsum[35-37]

\section{Limitations}\label{sec:disc-limitations}

\lipsum[38-40]

\section{Comparison with prior work}\label{sec:disc-comparison}

\lipsum[41-42]


\chapter{Further Analysis}\label{ch6}
% =============================================================================
% CHAPTER 6: Further Analysis
% =============================================================================
% SPDX-FileCopyrightText: 2026 Frank Langbein <frank@langbein.org>
% SPDX-License-Identifier: LPPL-1.3c
%
% Optional chapter: sensitivity analysis, extensions, case studies, or additional
% experiments. Suggested sections: Analysis; Implications (or as needed).
% Replace with your content.
% =============================================================================

This chapter extends the analysis or presents supplementary results. It is organised into analysis and implications. Replace with your content.

\lipsum[43-44]

\section{Analysis}\label{sec:further-analysis}

\lipsum[45-47]

\section{Implications}\label{sec:further-implications}

\lipsum[48-50]


\chapter{Conclusions and Future Work}\label{ch7}
% =============================================================================
% CHAPTER 7: Conclusions and Future Work
% =============================================================================
% SPDX-FileCopyrightText: 2026 Frank Langbein <frank@langbein.org>
% SPDX-License-Identifier: LPPL-1.3c
%
% Suggested sections: Conclusions (summary of aims and results; no new material);
% Future work. Optional: list of contributions.
% Replace with your closing remarks.
% =============================================================================

This chapter summarises the conclusions of the thesis and outlines future work. Replace with your closing remarks.

\lipsum[51-54]

\section{Conclusions}\label{sec:conclusions}
% Summary of aims, key results, and main takeaways. Optional: list of contributions.

\lipsum[55-57]

\section{Future work}\label{sec:future}

\lipsum[58-60]


\backmatter

% -----------------------------------------------------------------------------
% Appendices (lettered A, B, ... in headers and TOC)
% -----------------------------------------------------------------------------
\appendix
\chapter{Supplementary Methods}\label{app:A}
% =============================================================================
% Appendix A: Supplementary methods (demo)
% =============================================================================
% SPDX-FileCopyrightText: 2026 Frank Langbein <frank@langbein.org>
% SPDX-License-Identifier: LPPL-1.3c
%
% Content that supports the main Methods chapter: protocols, derivations,
% additional formulae, or detailed procedures. Rename the chapter title in
% main.tex and replace with your own content.
% =============================================================================

This appendix provides supplementary methods, protocols, or derivations that support Chapter~\ref{ch3}. Replace with your own content.

\section{Additional derivations}

\lipsum[61-62]

\section{Notation}

\lipsum[67]


\chapter{Supplementary Figures and Tables}\label{app:B}
% =============================================================================
% Appendix B: Supplementary figures and tables (demo)
% =============================================================================
% SPDX-FileCopyrightText: 2026 Frank Langbein <frank@langbein.org>
% SPDX-License-Identifier: LPPL-1.3c
%
% Supplementary figures, tables, or data summaries. Rename the chapter title in
% main.tex and replace with your own content.
% =============================================================================

This appendix contains supplementary figures and tables. Replace with your own content.

\section{Additional tables}

\lipsum[63-64]

\section{Supplementary figures}

\lipsum[65-66]


% Bibliography (always last): \bib{bibfile} uses natbib (numeric or harvard per class option)
% and inserts the License chapter if \license was set.
\bib{bibliography}

\end{document}
